\providecommand{\vf}[1]{\textcolor{blue}{#1}}
\providecommand{\hit}[1]{\textcolor[RGB]{0,115,90}{#1}}
\providecommand{\miss}[1]{\textcolor[RGB]{196,78,0}{\textbf{#1}}}

\chapter{Weakly Supervised Attack with a Small Leak}\label{ch:attack}

Chapter~\ref{ch:hub} trained the translation hub with full access to the hidden states of its three models.
This chapter adds a target model and connects it to the frozen hub through an affine map.
The frozen projectors and decoders then reconstruct the target model's prompts from the output of the map.

Two sources of information train the map.
First, a small leak, a set of prompts with both their text and the target model's hidden states, supervises the map directly.
Second, a matching term aligns unpaired data, hidden states of the target model and hub outputs with no correspondence between them.
The chapter tests whether the matching term improves reconstruction over the leak alone.
With a leak of \vf{90} rows, about ninety words, the matching term raises mean word-F1 on \vf{4{,}500} held-out NQ prompts from \vf{0.35} to \vf{0.77} over three seeds. % src: results/v7_2/main/w46-confirmation-s{1,2,3}.json (0.3544, 0.7659)

%% ------------------------------------------------------------------
\section{Threat Model}\label{sec:attack:threat}

The attacker is the honest-but-curious server of a split inference system and receives the hidden states that a client transmits to the cloud.
The attacker aims to reconstruct a prompt from its sequence of hidden states.
For that prompt the attacker holds neither the text nor the token identities and cannot query the target model.
In the experiments, the target model is Pythia-1B, observed at block \vf{12} (width \vf{2048}) with one vector per token. % src: configs/v7.yaml (unknowns.pythia: d3_block 12, width 2048)

The attacker holds four resources.
\begin{itemize}
\item \textbf{The hub.} The frozen hub of Chapter~\ref{ch:hub} with its source adapters, its projectors, and its three decoders (Qwen, OLMo, SmolLM).
\item \textbf{Two unpaired pools.} The first pool holds hidden states of the target model for prompts of the domain under attack, without their text.
The second pool holds hub outputs of text from the same domain, computed by OLMo and its source adapter.
No prompt of one pool is matched to a prompt of the other.
\item \textbf{A leak.} A small set of prompts for which the attacker holds both the text and the target model's hidden states.
\item \textbf{The tokenizer of the target model}, applied to the leaked prompts only.
\end{itemize}

\paragraph{Rows.} OLMo turns the leak into paired data.
For each leaked prompt, the attacker runs OLMo on the text and obtains the hub outputs of that prompt. % src: src/v6_4/pairs.py (build_pairs)
The target model and OLMo split the text into different tokens, so the attacker groups their tokens by shared character boundaries, as in the token groups of Section~\ref{sec:hub:training}.
Each group forms a row, a pair of vectors for one short span of text.
The pair holds the average of the target model's hidden states and the average of the hub outputs, both weighted by character coverage and not normalized to unit length. % src: src/v6_4/pairs.py
A row covers about one word, so an NQ prompt of nine words yields about nine rows. % src: results/v7_2/main/mixed-oct_w5b-A.json anchor_counts (10 prompts, 92 rows)
The affine map is fitted on rows, so the size of a leak is counted in rows.

%% ------------------------------------------------------------------
\section{Method}\label{sec:attack:method}

The attack trains a single affine map.
The map sends each hidden state $s$ of the target model to a point $\hat m$ of the hub,
\begin{equation}
\hat m = W(s - \mu) + b ,
\label{eq:attack:map}
\end{equation}
where $\mu$ is the mean hidden state of the unpaired pool, $W$ is a matrix from \vf{2048} to \vf{512} dimensions, and $b$ is a bias. % src: src/v7/adapter.py; src/v7_2/preprocess.py (centering on pool U); configs/v7.yaml (width 2048, hub_dim 512)
The frozen projectors and decoders of the hub then reconstruct the prompt from the mapped states.
Like the hub, the map acts on one hidden state at a time (Section~\ref{sec:hub:design}).

\paragraph{Objective.} The objective combines a supervised term on the leak with a matching term on the unpaired pools,
\begin{equation}
\mathcal{L}(\alpha) = \alpha\,\big(\mathcal{L}_{\text{sup}} + \lambda\,\mathcal{R}\big) + (1-\alpha)\,\mathcal{L}_{\text{match}} ,
\label{eq:attack:loss}
\end{equation}
where the supervision fraction $\alpha$ weights the two terms and does not depend on the size of the leak. % src: results/v7_2/main/mixed-v7_9-s5.json (definition)
At $\alpha = 1$ only the supervised term and its ridge penalty remain, and the result is the supervised fit.
At $0 < \alpha < 1$ both terms act, and the result is the hybrid fit.
At $\alpha = 0$ only the matching term remains, and the result is a blind map.

\paragraph{Supervised term.} $\mathcal{L}_{\text{sup}}$ is the mean squared distance between the mapped hidden state and the hub output of the same row, divided by $\rho_H^2$.
The constant $\rho_H$ is the root-mean-square distance of the pooled hub outputs from their mean, and the division makes the term independent of the scale of the hub. % src: src/v7_2/context.py (rho_H); src/v7/loss.py
The ridge penalty $\mathcal{R}$ is the squared Frobenius norm of $W$ divided by $\rho_H^2$, with $\lambda = 1$. % src: configs/v7_9_s4.yaml (mixed.ridge_lambdas); src/v7_2/mixed.py

\paragraph{Matching term.} $\mathcal{L}_{\text{match}}$ aligns a batch of $B$ mapped hidden states $\hat m_1 \ldots \hat m_B$ from the target model's pool with a batch of $B$ hub outputs $h_1 \ldots h_B$ from the hub's pool.
No element of the first batch is paired with an element of the second.
The term scores a round trip, from a mapped state to the hub batch and back to the mapped states.
The score of a pair is its negative squared distance, $\sigma(u, v) = -\lVert u - v\rVert^2 / \rho_H^2$.
On the way out, each mapped state receives a weighted average of the hub batch, with weights that favor the nearest hub outputs,
\begin{equation}
\bar h_i = \sum_{j} p_{ij}\, h_j ,
\qquad
p_{ij} = \frac{\exp(\sigma(\hat m_i, h_j)/\tau)}{\sum_{k}\exp(\sigma(\hat m_i, h_k)/\tau)} .
\label{eq:attack:match}
\end{equation}
On the way back, the weighted average $\bar h_i$ must identify the mapped state it came from among the $B$ mapped states of the batch,
\begin{equation}
\ell_i = -\log\frac{\exp(\sigma(\bar h_i, \hat m_i)/\tau)}{\sum_{k}\exp(\sigma(\bar h_i, \hat m_k)/\tau)} .
\label{eq:attack:return}
\end{equation}
$\mathcal{L}_{\text{match}}$ averages $\ell_i$ over the batch and adds the same loss with the two batches exchanged, each with weight one half. % src: src/v7_2/abc_sqdist.py (abc_sqdist_loss)
The loss is low when each weighted average distinguishes its own mapped state from the other mapped states of the batch.
The term requires no correspondence between the hidden states of the target model and the hub outputs.

\paragraph{Temperature.} The temperature $\tau$ sets the sharpness of both steps of the round trip, the weights $p_{ij}$ and the identification of Equation~\ref{eq:attack:return}.
The temperature is expressed as a multiple of a scale $d$ measured on the hub alone.
For each hub output, the gap separates its highest and second-highest scores against the other hub outputs of its batch.
The scale $d$ is the median gap over \vf{100} batches of \vf{512} hub outputs. % src: src/v7_2/gap.py; results/*/v7_2/main/mixed-v7_9-s0.json gap
The scale is \vf{0.0576} on NQ, \vf{0.0648} on legal128w, and \vf{0.0738} on legal10w. % src: results/v7_2/main/mixed-v7_9-s0.json gap.sqdist.median (0.05761); results/legal128w/v7_2/main/mixed-v7_9-s0.json (0.06485); results/legal10w/v7_2/main/mixed-v7_9-s0.json (0.07380)

\paragraph{Training.} Every fit runs \vf{3{,}000} AdamW updates, each on \vf{512} rows of the leak and $B$ elements from each unpaired pool. % src: configs/v7_9_s4.yaml (training); src/v7/training.py (lr_multiplier); results/v7_2/main/mixed-v7_9-s5.json setting (paired_batch_rows 512)
The map after the last update is evaluated, without checkpoint selection.

\paragraph{Hypothesis H2: matching improves reconstruction from a small leak.}\label{sec:attack:hypotheses}
With a leak of \vf{90} rows, word-F1 is higher under the hybrid fit than under the supervised fit.
The criterion, written before the run, requires that on NQ and on legal128w the gain in word-F1 be at least \vf{0.10}, with the lower bound of its \vf{95} percent interval above \vf{0.05}. % src: w46/plan.md (prediction P1)
H2 concerns word-F1, the word overlap between a reconstruction and the true prompt.
Exact match and decoded examples are reported beside it to separate word overlap from full recovery.

%% ------------------------------------------------------------------
\section{Experimental Setup}\label{sec:attack:experiments}

\paragraph{Datasets and pools.} The experiments run on the three datasets of Chapter~\ref{ch:hub}, NQ, legal128w, and legal10w, each with its own frozen hub.
The main comparison and the leak-size comparison run on all three datasets, and the grid search and the other supporting experiments run on NQ, on legal128w, or on both.
The \vf{40{,}000} training prompts of a dataset are divided once, at random, into two halves of \vf{20{,}000}.
One half supplies the target model's unpaired pool as hidden states, and the other half supplies the hub's unpaired pool as hub outputs of OLMo. 
A leak draws its prompts at random from all \vf{40{,}000} training prompts. 

\paragraph{The 90-row leak.} In the main experiment, the attacker obtains a leak of \vf{90} rows, the first nine rows of each of ten prompts.
The leak exposes the datasets to different degrees.
On NQ, a prompt averages nine words, so nine rows expose almost the entire prompt.
On legal128w, a prompt is a window of about \vf{68} words, or \vf{86} rows, cut from a longer document, so the same nine rows expose only its opening words. % src: artifacts/legal128w/v3/manifests/main-splits.jsonl window_words, validation split (mean 67.6); results/legal128w/v7_2/main/v8_6-summary.json rows[*].pairs (8612 rows for 100 prompts)
On legal10w, a prompt holds about \vf{13} rows, so nine rows expose most of it. % src: results/legal10w/v7_2/main/v8_6-summary.json rows[*] (1, 3, 5, 10, 30, 100 prompts: 12, 34, 55, 125, 390, 1305 rows; E +0.0055, +0.2723, +0.300, +0.332, +0.220, +0.051; job 1595870, seed 1)
The attacker's knowledge is therefore fixed at \vf{90} rows on every dataset, while the exposed fraction of each leaked prompt differs.

\paragraph{Score.} Evaluation follows Section~\ref{sec:hub:eval}.
Each of the three decoders generates greedily from the projected output of the map, up to the length of the soft prompt plus \vf{32} tokens in the main comparison.
The score is word-F1, averaged over prompts and over the three decoders, and word-F1 without a decoder's name denotes that three-decoder mean. 

\paragraph{Control and hub reference.} The mismatched-prompt control gives each decoder the mapped states of another prompt and scores the output against the original prompt.
The hub reference is the word-F1 of the three decoders on the hub outputs of OLMo, with no map involved, \vf{0.946} on NQ and \vf{0.982} on each legal dataset. % src: results/{legal128w,legal10w}/v7_2/main/hub-dev-eval-m32.json olmo_source_mean.hub (0.9821, 0.9821), 500 selection prompts; NQ 0.946 untraced, results/v7_2/main/hub-dev-eval-m32.json gives 0.9468

\paragraph{Settings.} The hybrid fit has three settings, the supervision fraction $\alpha$, the temperature $\tau$, and the batch size $B$ of the matching term.
The settings are chosen on \vf{500} validation prompts, the selection prompts, and the main comparison is scored on the other \vf{4{,}500}, the held-out prompts. % src: v7_run.py (subset_seed 20260903, dev_prompts 500); src/v7_2/context.py; results/v7_2/main/eval/confirmation/prompt_ids.json
The main comparison and the leak-size comparison fix $\alpha = 0.9$, $\tau = 2d$, and $B = 8{,}192$ on every dataset, and each other supporting experiment states its setting with its result.

%% ------------------------------------------------------------------
\section{Matching Improves Reconstruction from a Small Leak}\label{sec:attack:hybrid}

One setting of the hybrid fit, selected on NQ and legal128w, serves all three datasets, and at that setting matching improves word-F1 over the supervised fit on the same leak.

\subsection{Supervision fraction and temperature}\label{sec:attack:grid}

The gain from matching depends on the supervision fraction and the temperature, so a grid search over both precedes the main comparison (Figure~\ref{fig:attack:grid}).
The grid runs on the \vf{90}-row leak, with seed \vf{1} and the \vf{500} selection prompts of each dataset. % src: results/v7_2/main/mixed-w46-grid.json; results/legal128w/v7_2/main/mixed-w46-grid.json; w46-selection.json
The selection rule, written before the grid was run, takes the cell with the highest mean word-F1 over the two datasets. % src: w46/plan.md; w46-selection.json rule
That cell is $\alpha = 0.9$, $\tau = 2d$, with a two-dataset mean of \vf{0.534}. % src: w46-selection.json best (0.5336)
The cell is the best on NQ (\vf{0.768}) and close to the best on legal128w, \vf{0.299} against \vf{0.314} at $\alpha = 0.8$, $\tau = 2d$. % src: mixed-w46-grid.json seeds.1.cells (0.7679; legal128w 0.2993, 0.3144)
On legal10w, a dataset outside the selection rule, the same cell is again the best, at \vf{0.604}. % src: results/legal10w/v7_2/main/mixed-w46-grid.json seeds.1.cells (a0.9, t0.1476: 0.6040; next a0.8, t0.1476: 0.580; supervised 0.2251); job 1600092

All three datasets show the same dependence on temperature.
For $\alpha \le 0.95$, word-F1 peaks at $\tau = 2d$ and falls on either side, and several cells at a low temperature score below the supervised fit.
The matching term therefore helps only within a range of temperatures.
One setting serves all three datasets, although finding it required prompts with known text.

\begin{figure}[ht]
\centering
\includegraphics[width=\textwidth]{ch4/grid_search}
\caption{One setting, $\alpha = 0.9$ and $\tau = 2d$ (outlined), has the highest mean word-F1 over NQ and legal128w, the two datasets of the selection rule, and is the best cell on legal10w.
Word-F1 of the hybrid fit by supervision fraction (rows) and temperature (columns), on the \vf{90}-row leak, seed \vf{1}, \vf{500} selection prompts per dataset, $B = 8{,}192$.
The titles give the supervised fit on the same leak.}
\label{fig:attack:grid}
\end{figure}

\subsection{The main comparison}\label{sec:attack:main}

Unpaired matching more than doubles reconstruction word-F1 on all three datasets and meets the criterion of H2 on NQ and legal128w, the two datasets that the criterion names (Table~\ref{tab:attack:main}, Figure~\ref{fig:attack:standard}).
On the \vf{4{,}500} held-out prompts, with a leak of \vf{90} rows, word-F1 rises from \vf{0.354} under the supervised fit to \vf{0.766} under the hybrid fit on NQ, and from \vf{0.102} to \vf{0.321} on legal128w. % src: results/v7_2/main/w46-confirmation-s{1,2,3}.json; results/legal128w/v7_2/main/w46-confirmation-s{1,2,3}.json (means of three seeds)
The gain is \vf{0.411} on NQ, with an interval of \vf{0.408} to \vf{0.415}, and \vf{0.219} on legal128w, with an interval of \vf{0.215} to \vf{0.222}. % src: computed from results/{v7_2,legal128w/v7_2}/main/mixed-v7_9/confirmation/*/decoding.jsonl (per-prompt word-F1 averaged over three decoders and three seeds; paired bootstrap over 4,500 prompts, 2,000 resamples, seed 20261113)
Both gains exceed the required \vf{0.10}, and both lower bounds exceed the required \vf{0.05}.
The gain appears in every seed, at no less than \vf{0.402} on NQ and \vf{0.201} on legal128w. % src: same files, hybrid minus supervised per seed (0.4018, 0.4189, 0.4135; 0.2010, 0.2444, 0.2112)
On legal10w, a dataset added after the criterion was written, word-F1 rises from \vf{0.192} to \vf{0.599}, a gain of \vf{0.407} with an interval of \vf{0.403} to \vf{0.412}. % src: results/legal10w/v7_2/main/w46-confirmation-s{1,2,3}.json (supervised 0.1917, hybrid 0.5991; per seed 0.2229, 0.1989, 0.1532 and 0.6028, 0.5952, 0.5993); interval computed from results/legal10w/v7_2/main/mixed-v7_9/confirmation/*/decoding.jsonl (per-prompt word-F1 averaged over three decoders and three seeds; paired bootstrap over 4,500 prompts, 2,000 resamples, seed 20261113: +0.4074 [0.4033, 0.4118])
The legal10w gain is at least \vf{0.380} in every seed. % src: same files, hybrid minus supervised per seed (0.3799, 0.3963, 0.4461)
Both fits receive the same leak, so the gain comes from the two unpaired pools.

\begin{table}[H]
\caption{With \vf{90} leaked rows, the hybrid fit exceeds the supervised fit by \vf{0.41} word-F1 on NQ, by \vf{0.22} on legal128w, and by \vf{0.41} on legal10w.
Word-F1 averaged over three decoders and three seeds, Pythia as the target model, $\alpha = 0.9$, $\tau = 2d$, $B = 8{,}192$.
The held-out prompts selected no setting of the attack; each interval is a \vf{95} percent interval from a paired bootstrap over prompts.}
\label{tab:attack:main}
\centering
\small\setlength{\tabcolsep}{4pt}
\begin{tabular}{llccc}
\toprule
Dataset & Prompts & Supervised fit & Hybrid fit & Gain \\
\midrule
NQ & \vf{4{,}500} held-out & \vf{0.354} & \vf{0.766} & \vf{$+0.411$} [\vf{0.408}, \vf{0.415}] \\ % src: results/v7_2/main/w46-confirmation-s{1,2,3}.json cells.*.macro_f1 (0.3544, 0.7659)
NQ & \vf{500} selection & \vf{0.352} & \vf{0.762} & \vf{$+0.410$} [\vf{0.401}, \vf{0.420}] \\ % src: results/v7_2/main/mixed-w46-std.json seed_summary.budgets.90 (0.3519, 0.7623, 0.4104 [0.4005, 0.4196])
legal128w & \vf{4{,}500} held-out & \vf{0.102} & \vf{0.321} & \vf{$+0.219$} [\vf{0.215}, \vf{0.222}] \\ % src: results/legal128w/v7_2/main/w46-confirmation-s{1,2,3}.json (0.1021, 0.3209)
legal128w & \vf{500} selection & \vf{0.102} & \vf{0.320} & \vf{$+0.218$} [\vf{0.208}, \vf{0.227}] \\ % src: results/legal128w/v7_2/main/mixed-w46-std.json seed_summary.budgets.90 (0.1021, 0.3196, 0.2175 [0.2076, 0.2268])
legal10w & \vf{4{,}500} held-out & \vf{0.192} & \vf{0.599} & \vf{$+0.407$} [\vf{0.403}, \vf{0.412}] \\ % src: results/legal10w/v7_2/main/w46-confirmation-s{1,2,3}.json (supervised 0.1917, hybrid 0.5991; per seed 0.2229, 0.1989, 0.1532 and 0.6028, 0.5952, 0.5993); interval computed from results/legal10w/v7_2/main/mixed-v7_9/confirmation/*/decoding.jsonl (per-prompt word-F1 averaged over three decoders and three seeds; paired bootstrap over 4,500 prompts, 2,000 resamples, seed 20261113: +0.4074 [0.4033, 0.4118])
legal10w & \vf{500} selection & \vf{0.190} & \vf{0.602} & \vf{$+0.412$} [\vf{0.399}, \vf{0.425}] \\ % src: results/legal10w/v7_2/main/mixed-w46-std.json seeds.*.cells.*.final.macro_f1 (supervised 0.2251, 0.1958, 0.1489; hybrid 0.6040, 0.5947, 0.6060); seed_summary.budgets.90 (0.4116 [0.3990, 0.4251])
\bottomrule
\end{tabular}
\end{table}

\begin{figure}[ht]
\centering
\includegraphics[width=0.75\textwidth]{ch4/standard_cell}
\caption{The hybrid fit exceeds the supervised fit in every seed on all three datasets.
Word-F1 averaged over three decoders on the \vf{4{,}500} held-out prompts, with a leak of \vf{90} rows; bars are means over three seeds, dots the seeds, and each bracket the gain of the hybrid fit.}
\label{fig:attack:standard}
\end{figure}

\paragraph{Control and reference.} The hybrid scores depend on the prompt's own hidden states and remain below the hub reference.
On the selection prompts, the hybrid fit scores \vf{0.762} against its mismatched-prompt control of \vf{0.110} on NQ, \vf{0.320} against \vf{0.073} on legal128w, and \vf{0.602} against \vf{0.069} on legal10w. % src: mixed-w46-std.json seeds.*.cells.*.wrong_memory.macro_f1 (means 0.1104, 0.0733; legal10w 0.0690 from 0.0739, 0.0612, 0.0720)
The hub references are \vf{0.946} on NQ and \vf{0.982} on each legal dataset.
On legal128w the supervised fit, at \vf{0.102}, is close to its control of \vf{0.070}, so \vf{90} rows alone yield little beyond the control. % src: same file (0.0698)

\paragraph{Exact match and decoders.} Exact match distinguishes the word-overlap gain from full-prompt recovery.
Under the hybrid fit, exact match is \vf{0.21} on NQ and rounds to zero on legal128w. % src: computed from results/{v7_2,legal128w/v7_2}/main/mixed-v7_9/confirmation/*/decoding.jsonl, field exact_match, three decoders and three seeds (Corpus A hybrid 0.2121: SmolLM 0.4153, OLMo 0.1676, Qwen 0.0534; supervised 0.0001; legal128w hybrid: SmolLM 0.0001, OLMo 0, Qwen 0)
The three decoders also differ.
The hybrid fit scores \vf{0.886} through SmolLM, \vf{0.762} through OLMo, and \vf{0.650} through Qwen on NQ, and \vf{0.430}, \vf{0.160}, and \vf{0.372} on legal128w. % src: w46-confirmation-s{1,2,3}.json cells.*.per_decoder, means of three seeds (0.8859, 0.7622, 0.6495; 0.4303, 0.1602, 0.3722)
On legal10w the three scores are \vf{0.655}, \vf{0.587}, and \vf{0.555}. % src: results/legal10w/v7_2/main/w46-confirmation-s{1,2,3}.json cells.*.per_decoder, means of three seeds (SmolLM 0.6554, OLMo 0.5866, Qwen 0.5553)
Section~\ref{sec:attack:rewriting} draws on that difference.

\paragraph{Decoded examples.} The hybrid fit recovers content words that the supervised fit loses, as three held-out NQ prompts through the SmolLM decoder illustrate (Table~\ref{tab:attack:examples}).
At a word-F1 of \vf{0.80} the hybrid fit changes one name, and at \vf{0.40} it keeps the question frame and loses the title.
The supervised fit on the same leak keeps the question word and loses the content of all three prompts.

\begin{table}[H]
\caption{With the same \vf{90} rows, the hybrid fit reconstructs content words that the supervised fit loses.
Reconstructions through the SmolLM decoder for three held-out NQ prompts, chosen at high, intermediate, and low word-F1 of the hybrid fit (seed \vf{1}).
Decoded words that match the true prompt are green; the others are orange and bold.}
\label{tab:attack:examples}
\centering
\small
\begin{tabular}{@{}p{0.27\textwidth}p{0.68\textwidth}@{}}
\toprule
True prompt & when did kindergarten start in the united states \\
Hybrid fit, word-F1 \vf{1.00} & \hit{when did kindergarten start in the united states} \\ % src: results/v7_2/main/mixed-v7_9/confirmation/b90-a0.9-sq-m0-t0.11522-r1-B8192-L90-span9/decoding.jsonl, nq-00000406, smollm
Supervised fit, word-F1 \vf{0.47} & \hit{when did} \miss{semay faqum fell} \hit{in the} \miss{nile greece} \\ % src: results/v7_2/main/mixed-v7_9/confirmation/b90-a1-r1-L90-span9/decoding.jsonl, nq-00000406, smollm
\midrule
True prompt & who played sweet daddy williams on good times \\
Hybrid fit, word-F1 \vf{0.80} & \hit{who played} \miss{sweetheart} \hit{williams on good times} \\ % src: same files, nq-00001253
Supervised fit, word-F1 \vf{0.13} & \hit{who} \miss{sang sing choi will ton ghiam xo} \\
\midrule
True prompt & who played robin on once upon a time \\
Hybrid fit, word-F1 \vf{0.40} & \hit{who played} \miss{roon} \hit{on} \miss{one after another} \\ % src: same files, nq-00020139
Supervised fit, word-F1 \vf{0.12} & \hit{who} \miss{won homer in many times until he won} \\
\bottomrule
\end{tabular}
\end{table}

%% ------------------------------------------------------------------
\section{Where the Gain Holds}\label{sec:attack:ablations}\label{sec:attack:negative}

Five comparisons delimit the main result.
They vary the presence of the leak, its size, its composition, the components of the matching term, and the target model.

\subsection{Matching without a leak}\label{sec:attack:blind}

Matching alone, without any leaked row, yields little word-F1 beyond the mismatched-prompt controls.
On NQ, the blind maps tested score at most \vf{0.153} word-F1, against a hub reference of \vf{0.946}, and none exceeds its mismatched-prompt control by more than \vf{0.025} (interval \vf{0.019} to \vf{0.030}). % src: results/v7_2/main/gw.json (highest blind word-F1 0.1528, control 0.1414); results/v7_2/main/neighbourhood.json seeds.1.comparison.factorial (0.0801, 0.0555, +0.0247 [0.0187, 0.0304]); also results/v7/main-different_text-pythia/results.json and results/v7_2/main/mixed-v7_9-s1.json (alpha 0 cells)
The blind-map gains remain small compared with the gains from matching on a small leak.
The result covers \vf{22} blind maps from two objectives, one target model, and one dataset, and it does not show that no unpaired method can connect a target model. % src: 9 + 1 + 6 + 6 maps in the four files above

\subsection{Size of the leak}\label{sec:attack:budget}

Matching helps most on leaks of \vf{27} to \vf{270} rows (Figure~\ref{fig:attack:budget}).
The comparison repeats the main comparison at six leak sizes, the first nine rows of \vf{1}, \vf{3}, \vf{5}, \vf{10}, \vf{30}, and \vf{100} prompts, or \vf{9} to \vf{900} rows. % src: configs/w46_{nq,legal128w,legal10w}_std_L{9,27,45,270,900}.yaml; src/v7_2/anchors_linear.py span_start_pairs
The setting stays at $\alpha = 0.9$, $\tau = 2d$, and $B = 8{,}192$, with one seed on the selection prompts.
At \vf{90} rows the curve passes through the main comparison, where seed \vf{1} gains \vf{0.401} on NQ, \vf{0.206} on legal128w, and \vf{0.379} on legal10w. % src: results/{v7_2,legal128w/v7_2,legal10w/v7_2}/main/mixed-w46-std.json seeds.1.cells (hybrid 0.7679, 0.2993, 0.6040; supervised 0.3671, 0.0938, 0.2251)
On NQ, the gain is \vf{0.354} at \vf{27} rows, peaks at \vf{0.418} at \vf{45} rows, and falls to \vf{0.270} at \vf{270} rows and \vf{0.089} at \vf{900} rows. % src: results/v7_2/main/mixed-w46-std-L{27,45,270,900}.json seeds.1.cells (hybrid 0.6093, 0.7217, 0.7956, 0.8323; supervised 0.2549, 0.3037, 0.5255, 0.7434); job 1601532
On legal10w, the gain is \vf{0.243} at \vf{27} rows, peaks at \vf{0.379} at \vf{90} rows, and falls to \vf{0.314} at \vf{270} rows and \vf{0.089} at \vf{900} rows. % src: results/legal10w/v7_2/main/mixed-w46-std-L{27,45,270,900}.json (hybrid 0.3799, 0.5163, 0.6880, 0.7437; supervised 0.1365, 0.1686, 0.3737, 0.6548; gain at 45 rows 0.348); job 1601534
On legal128w, the gain is \vf{0.090} at \vf{27} rows and \vf{0.198} at \vf{45} rows, peaks at \vf{0.206} at \vf{90} rows, and falls to \vf{0.173} at \vf{270} rows and \vf{0.018} at \vf{900} rows. % src: results/legal128w/v7_2/main/mixed-w46-std-L{27,45,270,900}.json seeds.1.cells (hybrid 0.1483, 0.2571, 0.3878, 0.4960; supervised 0.0586, 0.0593, 0.2144, 0.4783); job 1601533
A leak of one prompt is too small for either fit.
At \vf{9} rows the hybrid fit scores \vf{0.195} on NQ, a gain of \vf{0.036}, and on the two legal datasets it stays at the level of its mismatched-prompt control, at \vf{0.021} and \vf{0.061}. % src: mixed-w46-std-L9.json of each dataset (NQ hybrid 0.1951, supervised 0.1590; legal128w hybrid 0.0207, supervised 0.0445, control 0.023; legal10w hybrid 0.0613, supervised 0.0814, control 0.027)
At \vf{900} rows the supervised fit has risen to \vf{0.743} on NQ and \vf{0.655} on legal10w, and little remains for matching to add.
The setting was selected at \vf{90} rows, so the gain at \vf{900} rows describes that setting and not the best hybrid fit for a leak of that size.

\begin{figure}[ht]
\centering
\includegraphics[width=\textwidth]{ch4/leak_size_three}
\caption{The gain from matching is largest at leaks of \vf{45} to \vf{90} rows and small at \vf{9} and \vf{900} rows.
Word-F1 averaged over three decoders on the \vf{500} selection prompts, seed \vf{1}, for the supervised fit and the hybrid fit of the main comparison ($\alpha = 0.9$, $\tau = 2d$, $B = 8{,}192$); each leak holds the first nine rows of \vf{1} to \vf{100} prompts, and the number above each point is the gain.
The dotted line is the hub reference.}
\label{fig:attack:budget}
\end{figure}

\todo{Appendix: this part stays in the chapter for now and moves to an appendix.}
A larger leak also moves the best setting of the grid of Section~\ref{sec:attack:grid} (Figure~\ref{fig:attack:grid-size}).
With one row from each of \vf{90} legal128w prompts, the best cell is $\alpha = 0.8$, $\tau = 2d$, at a word-F1 of \vf{0.341} against \vf{0.108} for the supervised fit. % src: results/legal128w/v7_2/main/mixed-oct_w5-grid-L90.json selection.winners.90 (0.3411); cells b90/a1/r1/L90-random (0.1084)
With one row from each of \vf{864} prompts, the best cell is $\alpha = 0.95$, $\tau = d$, at \vf{0.659} against \vf{0.548}. % src: results/legal128w/v7_2/main/mixed-oct_w5-grid-L864.json selection.winners.864 (0.6594); cells b864/a1/r1/L864-random (0.5480)
The larger leak is also less sensitive to the setting, since every cell with $\alpha \ge 0.9$ and $\tau \ge 0.5d$ lies between \vf{0.46} and \vf{0.66}. % src: same file, seeds.1.cells (0.4620 to 0.6594)

\begin{figure}[ht]
\centering
\includegraphics[width=0.95\textwidth]{ch4/grid_leak_size}
\caption{A larger leak moves the best cell toward a higher supervision fraction and a lower temperature (outlined), and flattens the grid.
Word-F1 of the hybrid fit by supervision fraction (rows) and temperature (columns) on legal128w, for leaks of one row from each of \vf{90} or \vf{864} prompts, seed \vf{1}, \vf{500} selection prompts, $B = 8{,}192$.
The titles give the supervised fit on the same leak.}
\label{fig:attack:grid-size}
\end{figure}

\subsection{Composition of the leak}\label{sec:attack:leak}

At a fixed number of rows, a leak spread over more prompts yields a higher word-F1 than a leak drawn from one or two prompts (Figure~\ref{fig:attack:shape}).
On legal128w, the hybrid fit scores \vf{0.321} with one row from each of \vf{90} prompts, \vf{0.304} with nine consecutive rows from each of ten prompts, and \vf{0.178} with \vf{90} rows from two prompts. % src: results/legal128w/v7_2/main/mixed-oct_w5-s1-L90.json (L90-random 0.3207); mixed-oct_w5-s1-L90-ext2.json (L90-span9 0.3039); mixed-oct_w5b-L-prefix.json (L90-prefix 0.1782)
The content of the rows matters as well.
A leak of stop words, or of one repeated word, leaves the hybrid fit at the level of the supervised fit, at \vf{0.058} against \vf{0.056} and \vf{0.054} against \vf{0.057}. % src: mixed-oct_w5-s1-L90.json cells *L90-stop (0.0582, 0.0555), *L90-one_word (0.0535, 0.0568)
The comparison rests on one seed and does not separate the number of prompts from the content of the particular prompts drawn.

At larger leaks the number of prompts no longer separates the two fits on legal128w.
Ten whole prompts supply \vf{856} rows, and the hybrid fit of the main comparison scores \vf{0.501} against \vf{0.487} for the supervised fit. % src: results/legal128w/v7_2/main/mixed-v8_7-s1.json seeds.1.cells b10/a0.9/sq/m0/t0.1297/r1/B8192 (0.5006), b10/a1/r1 (0.4869); budgets.10.anchor_counts.valid_aligned_pairs (856)
Nine rows from each of \vf{100} prompts supply \vf{900} rows, and the two fits score \vf{0.496} and \vf{0.478}. % src: results/legal128w/v7_2/main/mixed-w46-std-L900.json seeds.1.cells (0.4960, 0.4783)
The gain is below \vf{0.02} under both leaks, so at this size the fixed setting, and not the spread of the rows over prompts, limits the hybrid fit.
With \vf{30} and \vf{100} whole prompts, the supervised fit overtakes the hybrid fit, at \vf{0.731} against \vf{0.613} and \vf{0.813} against \vf{0.654}. % src: mixed-v8_7-s1.json, b30 and b100 cells (0.7307, 0.6131; 0.8133, 0.6544); 2,569 and 8,612 rows

\begin{figure}[ht]
\centering
\includegraphics[width=0.8\textwidth]{ch4/leak_shape}
\caption{With \vf{90} leaked rows on legal128w, the hybrid fit scores \vf{0.18} for a leak from two prompts and \vf{0.30} to \vf{0.32} for a leak spread over \vf{10} or \vf{90} prompts.
Word-F1 averaged over three decoders, seed \vf{1}, \vf{500} selection prompts, hybrid fit with $\alpha = 0.99$, $\tau = d$, $B = 8{,}192$.
Each bracket gives the gain of the hybrid fit.
The NQ leak takes whole prompts, so the two datasets differ in more than the domain.}
\label{fig:attack:shape}
\end{figure}

\subsection{Components of the matching term}\label{sec:attack:components}

\todo{Appendix: this part stays in the chapter for now and moves to an appendix.}
The matching term was selected in sequential steps on NQ, and four steps each raise word-F1 at a leak of about \vf{92} rows.
The squared-distance score adds \vf{0.043} over a cosine score, a supervision fraction of \vf{0.95} adds \vf{0.096} over one, a temperature of $2d$ adds \vf{0.110} over $d$, and a batch of \vf{2{,}048} adds \vf{0.092} over \vf{512}. % src: results/v7_2/main/mixed-v7_9-s{1,2,3}.json (score_effect, cell_minus_benchmark, axis_effect at budget 10), three seeds, 500 selection prompts

\subsection{A second target model}\label{sec:attack:llama}

The gain recurs for a second target model in the one available comparison.
Llama-3.2-1B, observed at block \vf{12} (width \vf{2048}), replaces Pythia on NQ with a leak of \vf{90} rows from ten prompts, one seed, and the selection prompts. % src: configs/v7.yaml (unknowns.llama: d3_block 12, width 2048)
Llama reaches a word-F1 of \vf{0.632} under the hybrid fit against \vf{0.370} under the supervised fit, and Pythia under the same setting reaches \vf{0.604} against \vf{0.352}. % src: results/nqllama/v7_2/main/mixed-oct_w5b-s1-prefix.json cells *L90-prefix (0.6324, 0.3696, +0.2628 [0.2514, 0.2742]); results/v7_2/main/mixed-oct_w5b-A-prefix.json (0.6042, 0.3516, +0.2526 [0.2426, 0.2640])

%% ------------------------------------------------------------------
\section{The Training Objective and the Decoder's Predictions}\label{sec:attack:teacher}

\todo{Appendix: this part stays in the chapter for now and moves to an appendix.}
Both terms of the objective in Equation~\ref{eq:attack:loss} pull mapped states toward hub outputs.
The mapped states then form the soft prompt of a frozen decoder, and the decoder reconstructs the prompt one token at a time.
The objective therefore measures agreement with hub outputs, while reconstruction depends on the decoder's predictions.

To examine how the decoder's predictions relate to the reconstruction score, the same decoder runs twice on each prompt.
The reference run receives a soft prompt from the hub outputs of OLMo, and the mapped-state run receives a soft prompt from the mapped states of the target model.
Both runs see the true previous tokens, a setting called teacher forcing, and give a probability to every possible next token.
A change in the predictions between the two runs therefore comes from replacing hub outputs with mapped states.
Both runs receive the true previous tokens, a setting called teacher forcing, and assign a probability to every candidate next token.
The comparison therefore shows how the decoder's predictions change when mapped states replace hub outputs.

Cross-entropy is the decoder's surprise at the true next token, the negative logarithm of the probability it assigns to that token.
A higher probability gives a lower cross-entropy.
The excess cross-entropy subtracts the cross-entropy of the reference run from that of the mapped-state run, so a positive excess means that the mapped states leave the decoder more surprised at the true tokens.

The KL divergence compares the two runs over all candidate tokens, not only the true one.
It measures how far the next-token distribution of the mapped-state run departs from that of the reference run, and it is zero when the two distributions match.
The excess cross-entropy thus measures the change in surprise at the true tokens, and the KL divergence measures the departure from the reference predictions over all tokens.

Both scores are summed over the tokens of each word and averaged over content words, across the selection prompts, the three decoders, and the seeds of the map.
The averages are reported in nats per content word. % src: teacher_forced_run.py (lines 62-68: -logp at the true token; KL(reference || mapped) = sum q (log q - log p)); teacher_forced_summary.py (lines 82-136: per-word sums, content words, pooling over prompts, decoders, seeds)

\paragraph{The matching term and the decoder.} On NQ the hybrid fit has the lower excess cross-entropy at every leak size, \vf{1.76} against \vf{7.61} nats at \vf{92} rows (Figure~\ref{fig:attack:ce}). % src: results/v7_2/main/framework/teacher_forced/summary.json arms.{b10-a1-r1, b10-a0.95-sq-m0-t0.1152-r1-B8192}.classes.content.excess_per_word (7.6127, 1.7640)
At that leak size the supervised fit, at \vf{7.61}, is worse than a decoder with no soft prompt at all, at \vf{6.08}. % src: same file, arms.no_memory (6.0751)
On legal128w the order reverses, and the hybrid fit has the higher excess cross-entropy at every leak size, \vf{2.82} against \vf{2.46} at \vf{856} rows. % src: results/legal128w/v7_2/main/framework/teacher_forced/summary.json (2.8195, 2.4597)
Its word-F1 at that leak size is nevertheless the higher one, \vf{0.523} against \vf{0.487}. % src: results/legal128w/v7_2/main/v8_6-summary.json rows (hybrid_B8192 0.5228, supervised 0.4869)
The KL divergence follows the excess cross-entropy within \vf{0.25} nats in every cell.
The matching term therefore raises word-F1 without consistently improving the decoder's next-token predictions.

\begin{figure}[ht]
\centering
\includegraphics[width=0.95\textwidth]{ch4/cross_entropy}
\caption{Under teacher forcing, the hybrid fit is below the supervised fit on NQ and above it on legal128w.
Excess cross-entropy in nats per content word against the number of leaked rows, \vf{500} selection prompts; the dotted line, labeled no memory, is a decoder with no soft prompt.
Hybrid fit with $\alpha = 0.95$, $\tau = 2d$, $B = 8{,}192$.}
\label{fig:attack:ce}
\end{figure}

\paragraph{Training the map on the decoder's predictions.}\label{sec:attack:decoderkl} A further step places the decoder inside the objective.
Starting from a hybrid fit on legal128w ($\alpha = 0.99$, $\tau = d$), the step trains the affine map through the frozen SmolLM decoder on the prompts of the leak, with the KL divergence defined above as its training loss. % src: oct_w4/plan.md; results/legal128w/v7_2/main/oct_w4-summary.json budgets.*.selected (kl/beta0)
On the \vf{4{,}500} held-out prompts, with one seed, the KL divergence and the cross-entropy of SmolLM both fall at every leak size (Table~\ref{tab:attack:decoderkl}).
The cross-entropy falls by \vf{0.19}, \vf{0.15}, and \vf{0.26} nats per content word at leaks of \vf{856}, \vf{2{,}569}, and \vf{8{,}612} rows.
Word-F1 rises by at most \vf{0.023} at the same leak sizes. % src: oct_w4-summary.json budgets.*.confirmation.gain_vs_start.kl/beta0.ce; confirmation.f1_gain_vs_start.kl/beta0.f1 (+0.0092 [0.0062, 0.0123]; +0.0111 [0.0090, 0.0134]; +0.0233 [0.0212, 0.0254])
%% BIG CHANGE: the author's reading (2026-11-15): the training loss falls, cross-entropy and KL divergence fall with it, and word-F1 gains little.
A lower training loss thus lowers the cross-entropy and the KL divergence but adds little word-F1.

\begin{table}[H]
\caption{Training the map through the decoder with the KL loss lowers cross-entropy and KL divergence by \vf{0.15} to \vf{0.26} nats and raises word-F1 by \vf{0.01} to \vf{0.02}.
SmolLM decoder, legal128w, seed \vf{1}, \vf{4{,}500} held-out prompts; cross-entropy and KL divergence in nats per content word under teacher forcing; each gain carries a \vf{95} percent interval from a paired bootstrap over prompts.}
\label{tab:attack:decoderkl}
\centering
\small\setlength{\tabcolsep}{4pt}
\begin{tabular}{lccc}
\toprule
Leak (rows) & Cross-entropy, before to after & KL divergence, before to after & Gain in word-F1 \\
\midrule
\vf{856} & \vf{1.48} to \vf{1.29} & \vf{1.45} to \vf{1.25} & \vf{$+0.009$} [\vf{0.006}, \vf{0.012}] \\ % src: results/legal128w/v7_2/main/oct_w4-summary.json budgets.10.confirmation.table.{start,kl/beta0} (1.4842, 1.2896; 1.4464, 1.2503); f1_gain_vs_start (+0.0092 [0.0062, 0.0123])
\vf{2{,}569} & \vf{1.02} to \vf{0.87} & \vf{1.00} to \vf{0.84} & \vf{$+0.011$} [\vf{0.009}, \vf{0.013}] \\ % src: same, budgets.30 (1.0236, 0.8709; 0.9966, 0.8432; +0.0111 [0.0090, 0.0134])
\vf{8{,}612} & \vf{0.77} to \vf{0.51} & \vf{0.75} to \vf{0.50} & \vf{$+0.023$} [\vf{0.021}, \vf{0.025}] \\ % src: same, budgets.100 (0.7703, 0.5131; 0.7520, 0.4962; +0.0233 [0.0212, 0.0254])
\bottomrule
\end{tabular}
\end{table}

%% ------------------------------------------------------------------
\section{Second-Round Rewriting}\label{sec:attack:rewriting}

\todo{Appendix: this part stays in the chapter for now and moves to an appendix.}
The three decoders yield reconstructions of differing quality (Section~\ref{sec:attack:main}), and a second round tests whether combining them improves word-F1 over their mean and over the best decoder.
A language model receives the three reconstructions and rewrites them into one text.
The experiment is exploratory, and no criterion preceded its runs.

\paragraph{Protocol.} The language model is Claude Haiku~\vf{5.5}.
It receives the three reconstructions, with the SmolLM reconstruction first, and returns one text. % src: results/v7_2/main/framework/attacker/attacker-b10-a0.95-sq-m0-t0.1152-r1-B8192-s1-haiku.json (model claude-haiku-5-5, system_B); attacker_run.py (LEVEL2_ORDER)
Its instruction reads ``You get three noisy copies of the same question. Copy 1 is the most reliable; start from it and apply majority rules where the copies disagree. Fix spelling, broken words and repeats. Reply with the text only, inside \texttt{<text></text>}.''
On the two legal datasets the instruction says ``SEC filing excerpt'' in place of ``question''.
The language model receives the reconstructions, the kind of text, and the instruction that names SmolLM as the most reliable copy, but no hidden state and no part of the true prompt.
The experiment covers \vf{450} of the \vf{500} selection prompts and one seed, and the instruction was selected on the other \vf{50}. % src: same files (prompts_reported 450); attacker_run.py (selection sample: first 50 prompt ids)

\paragraph{Result.} Rewriting exceeds the mean word-F1 of the three reconstructions in every tested condition, and its small advantage over the best single decoder carries no interval (Table~\ref{tab:attack:rewriting}).
On NQ at a leak of \vf{92} rows, word-F1 rises from a three-decoder mean of \vf{0.727} to \vf{0.861}, while SmolLM alone scores \vf{0.854}. % src: results/v7_2/main/framework/attacker/attacker-b10-a0.95-sq-m0-t0.1152-r1-B8192-s1-haiku.json; results/legal128w/v7_2/main/framework/attacker/attacker-b{10,100}-a0.95-sq-m0-t0.1297-r1-B8192-s1-haiku.json, table."hybrid rung B".word_f1; SmolLM alone recomputed from the per-prompt records (0.8539)
Across the four conditions the gain over the mean is \vf{0.12} to \vf{0.26}, with every paired interval above zero.
The gain over SmolLM alone is \vf{0.01} to \vf{0.02} on NQ and legal128w and \vf{0.06} on legal10w, where the three decoders score closest to one another. % src: results/legal10w/v7_2/main/framework/attacker/attacker-b90-a0.9-sq-m0-t0.1476-r1-B8192-L90-span9-s1-haiku.json table."hybrid rung B".word_f1 (0.6273, 0.7446); per-decoder scores recomputed from the endpoint decoding on the 450 prompts (SmolLM 0.6839, OLMo 0.6176, Qwen 0.5804)
Given one reconstruction, the language model changes word-F1 by at most \vf{0.023}, averaged over the three decoders. % src: same files, table."hybrid rung A".word_f1.gain (-0.0017, +0.0226, +0.0119)
Entity recall measures the share of the true prompt's named entities that appear in the output.
At leaks of \vf{92} rows on NQ and \vf{856} rows on legal128w, rewriting raises entity recall above the three-decoder mean, from \vf{0.50} to \vf{0.71} and from \vf{0.21} to \vf{0.37}. % src: same files, table."hybrid rung B".entity_recall (0.5000, 0.7116; 0.2106, 0.3743; 0.4399, 0.5735)
At \vf{8{,}612} rows on legal128w, entity recall rises from \vf{0.44} to \vf{0.57}.
On legal10w at \vf{90} rows, entity recall rises from \vf{0.17} to \vf{0.21}. % src: legal10w attacker file, table."hybrid rung B".entity_recall (0.1695, 0.2143, +0.045 [0.019, 0.074])

\begin{table}[H]
\caption{Rewriting exceeds the mean of the three decoders by \vf{0.12} to \vf{0.26} word-F1 and SmolLM alone by \vf{0.01} to \vf{0.06}.
Hybrid fits with $\tau = 2d$, $B = 8{,}192$, seed \vf{1}, \vf{450} selection prompts; $\alpha = 0.95$ on leaks of whole prompts for NQ and legal128w, and the fit of the main comparison ($\alpha = 0.9$, the \vf{90}-row leak) for legal10w; the gain is after rewriting minus the mean of three, with a \vf{95} percent interval from a paired bootstrap over prompts.
Word-F1 in this section is computed on lower-cased words without edge punctuation.}
\label{tab:attack:rewriting}
\centering
\small\setlength{\tabcolsep}{4pt}
\begin{tabular}{lcccc}
\toprule
Condition & Mean of three & SmolLM alone & After rewriting & Gain \\
\midrule
NQ, \vf{92} rows & \vf{0.727} & \vf{0.854} & \vf{0.861} & \vf{$+0.134$} [\vf{0.124}, \vf{0.145}] \\ % src: attacker-b10-a0.95-sq-m0-t0.1152-r1-B8192-s1-haiku.json table."hybrid rung B".word_f1 (0.7271, 0.8610, +0.1339 [0.1237, 0.1449]); SmolLM alone recomputed from the per-prompt records (0.8539)
legal128w, \vf{856} rows & \vf{0.539} & \vf{0.779} & \vf{0.802} & \vf{$+0.263$} [\vf{0.252}, \vf{0.273}] \\ % src: legal128w attacker-b10-...-haiku.json (0.5393, 0.8020, +0.2627 [0.2522, 0.2729]); SmolLM alone recomputed (0.7793)
legal128w, \vf{8{,}612} rows & \vf{0.772} & \vf{0.890} & \vf{0.896} & \vf{$+0.124$} [\vf{0.115}, \vf{0.135}] \\ % src: legal128w attacker-b100-...-haiku.json (0.7721, 0.8963, +0.1242 [0.1145, 0.1349]); SmolLM alone recomputed (0.8897)
legal10w, \vf{90} rows & \vf{0.627} & \vf{0.684} & \vf{0.745} & \vf{$+0.117$} [\vf{0.104}, \vf{0.130}] \\ % src: results/legal10w/v7_2/main/framework/attacker/attacker-b90-a0.9-sq-m0-t0.1476-r1-B8192-L90-span9-s1-haiku.json table."hybrid rung B".word_f1 (0.6273, 0.7446, +0.1173 [0.1041, 0.1303]); SmolLM alone recomputed (0.6839); run 2026-11-17, 4,500 calls
\bottomrule
\end{tabular}
\end{table}

\paragraph{What the result does not show.} The comparison does not separate the contribution of the two further reconstructions from the instruction that names SmolLM as the most reliable copy.
The rewriting by Claude Haiku also lacks a mismatched-prompt control.
A confirmatory test requires a paired interval for the gain over the best decoder and a mismatched-prompt control for the same protocol.

%% ------------------------------------------------------------------
\section{Discussion and Limitations}\label{sec:attack:discussion}

A target model outside hub training connects to the hub through one affine map, and unpaired matching improves the reconstruction obtained from a small supervised leak.
With \vf{90} leaked rows, word-F1 on held-out prompts rises from \vf{0.354} to \vf{0.766} on NQ, from \vf{0.102} to \vf{0.321} on legal128w, and from \vf{0.192} to \vf{0.599} on legal10w. % src: results/{v7_2,legal128w/v7_2,legal10w/v7_2}/main/w46-confirmation-s{1,2,3}.json
The gains on NQ and legal128w, the two datasets that H2 names, meet its criterion.

The attack trains only an affine map, and the hub and its decoders stay frozen.
An affine map preserves the linear structure of the hidden states it receives.
The gain from matching therefore suggests that the hidden states of the target model and the hub outputs share that structure up to an affine transformation.
The Platonic Representation Hypothesis predicts this similarity between models, here for a target model that took no part in hub training.

Three findings delimit the result.
The gain is largest at small leaks and shrinks as the leak grows, to \vf{0.076} at \vf{986} rows on NQ.
The gain holds only within a range of temperatures.
Without a leak, the tested blind maps exceed their mismatched-prompt controls by at most \vf{0.025} word-F1.

A lower decoder loss and a higher word-F1 do not improve together consistently.
The objective does not measure the decoder's predictions, and teacher-forced cross-entropy and KL divergence do.
Training the map on the KL divergence lowers the cross-entropy of SmolLM by \vf{0.15} to \vf{0.26} nats per content word on held-out legal128w prompts, with one seed, and adds little word-F1. % src: oct_w4-summary.json budgets.*.confirmation.gain_vs_start.kl/beta0.ce
A second-round language model merges the three decoders' reconstructions into one text at the level of the best decoder, \vf{0.12} to \vf{0.26} word-F1 above their mean.

\paragraph{What the reconstructions contain.} The results establish a gain in word overlap, not exact recovery, recovery of the facts of a prompt, or an amount of leaked information.
Exact match under the hybrid fit is \vf{0.21} on NQ and rounds to zero on legal128w.
In legal128w reconstructions from the leak-composition experiment (Section~\ref{sec:attack:leak}), the errors are fluent substitutions, e.g., ``registrant'' becomes ``insurer'', and dates, section numbers, and a company name are replaced by others of the same kind. % src: results/legal128w/v7_2/main/mixed-v7_9/runs/b90-a0.99-sq-m0-t0.06485-r1-B8192-L90-random-s1/endpoint/decoding.jsonl, SmolLM records
A reconstruction does not mark which of its words match the source.
At a leak of \vf{856} rows on legal128w, rewriting raises entity recall from \vf{0.21} to \vf{0.37}, and most named entities remain unrecovered.

\paragraph{Limitations.} Four limitations bound the claims.
\begin{itemize}
\item \textbf{Settings chosen with known text.} The supervision fraction and the temperature were selected by word-F1 on \vf{500} prompts whose text is known, and no run measures the cost of selecting them without such prompts.
\item \textbf{Supporting experiments.} Only the main comparison runs on held-out prompts with three seeds, and most supporting experiments rest on the \vf{500} selection prompts and one seed, at earlier settings.
\item \textbf{Same-domain pools.} Both unpaired pools and the leak come from the dataset under attack, and no run tests text from another domain or a smaller pool.
\item \textbf{Models and layer.} The results cover one block of two target models of about one billion parameters and three hub models, without a defense or a deployed system.
\end{itemize}

The experiments quantify reconstruction under the stated attacker resources, and an information-theoretic account remains future work (Chapter~\ref{ch:future}).
